% TOP_PROBLEM: {"id": 3000056, "title": "Minimum k-way cut in a hypergraph", "category_id": 2, "category": {"id": 2, "name": "combinatorics", "display_name": "Combinatorics", "description": "Counting problems, graph theory, discrete structures.", "slug": "combinatorics", "order_index": 2, "created_at": "2026-07-31T15:26:25.671Z"}, "status": "partially_solved"}
% FIRST_POSTED: null
% ENTRY_KIND: statement_only
% Imported from: batch09.tex; UnsolvedMath ID 3000056
\documentclass[11pt]{article}
\usepackage[margin=25mm]{geometry}
\usepackage{fontspec}
\setmainfont{Latin Modern Roman}
\usepackage{amsmath,amssymb,mathtools}
\usepackage{unicode-math}
\setmathfont{Latin Modern Math}
\usepackage{xurl,hyperref}
\hypersetup{colorlinks=true,urlcolor=blue}
\setcounter{secnumdepth}{0}
\setlength{\parindent}{0pt}
\setlength{\parskip}{5pt}
\setlength{\emergencystretch}{3em}
\newcommand{\sref}[2]{\hyperlink{source:#1:#2}{[S#2]}}
\newcommand{\cataloguescope}[1]{\paragraph{Recorded target.}#1}
\begin{document}
% UnsolvedMath ID 3000056; source catalogue code AMR-029-0056
\setcounter{section}{3000055}
\section{Minimum k-way cut in a hypergraph}\label{3000056}
{\small\sffamily UnsolvedMath ID 3000056 · Combinatorics}
\subsection{Definitions and mathematical statement}

\paragraph{Shared notation.}$\mathbb N=\{1,2,\ldots\}$, and $\mathbb Z,\mathbb Q,\mathbb R,\mathbb C$ denote the integers, rationals, reals and complex numbers. Any other conventions are specified locally.


Let $H=(V,\mathcal E)$ be a finite hypergraph, whose explicitly listed hyperedges may have arbitrary size. Each edge $e$ has a nonnegative binary-encoded rational capacity $c(e)$. Fix an integer $k\ge2$ with $k\le|V|$. A $k$-way cut is a partition $V=V_1\dot\cup\cdots\dot\cup V_k$ into nonempty parts; its cut set consists of the edges meeting at least two distinct parts. Its capacity is $\sum_{e\text{ crossing}}c(e)$, charging a crossing hyperedge once regardless of how many parts it meets. Equivalently, remove minimum-capacity hyperedges so that at least $k$ connected components remain.
\paragraph{Statement recorded in the supplied source.}
For every fixed $k$, can a minimum-capacity $k$-way hypergraph cut be found in time polynomial in the encoded input size, with no bound on edge rank? \sref{3000056}{2}
The historically registered question is answered affirmatively by Theorem 1.1 and the cost extension in Section 1.1 of \sref{3000056}{2}.
\subsection{Short English statement}
Can a cheapest k-way cut of an arbitrary capacitated hypergraph be computed in polynomial time when k is fixed?
\subsection{Sources}
\begin{description}
\item[\hypertarget{source:3000056:1}{S1}] ULAM.AI and the UnsolvedMath Contributors, \emph{UnsolvedMath: A Curated Collection of Open Mathematics Problems}, record 3000056 (AMR-029-0056). CC BY 4.0. \url{https://huggingface.co/datasets/ulamai/UnsolvedMath}.
\item[\hypertarget{source:3000056:2}{S2}] K. Chandrasekaran and C. Chekuri, Hypergraph k-cut for fixed k in deterministic polynomial time (2020), Introduction, formal Hypergraph-k-Cut definition, Theorem 1.1 and Section 1.1, pp.1--3. \url{https://arxiv.org/pdf/2009.12442}.
\end{description}
\label{end:3000056}
\end{document}
